\section{The singular contribution}\label{sec:tail}

We now prove that configurations with a positive point-pair intersection
make a vanishing contribution.  Throughout this section the ambient
symplectic space has dimension $2D$, where $D=2r$, and the point
Lagrangians $(L_i)_{i\in I}$ are independent and uniform.  We use the
unsigned normalized kernel $J$ from Section~\ref{sec:orbits}.  For an
activity set $B\subseteq\mathscr C$, put
\[
 F_j((L_i)_{i\in j})=
 \E_Y\prod_{i\in j:\,(i,j)\in B}J(L_i,Y).
\]
Center integration gives the product $\prod_jF_j$.  The bounds of
Section~\ref{sec:planted} are pointwise bounds on these factors, including
on exceptional configurations of the point variables.

For $e=\{i,k\}\in E$, set $h_e=\dim(L_i\cap L_k)$, and for $j\in J$ set
$c_j=\dim\bigcap_{i\in j}L_i$.  These dimensions form a \emph{profile}.
On a face we use opposite-pair indices $\ell=1,2,3$, so that
$h_\ell=c+s_\ell$, $d=D-c$, and
\[
 C_j=(D+1)c_j+\sum_{\ell=1}^3b(s_\ell),
 \qquad b(n)=\frac{n(n+1)}2.
\]
All constants in bounds of the form $O_D(q^x)$ may depend on $D$ and on
the fixed complex.  Bounds on an exponent that are stated to be uniform
in $D$ have a different role: they will allow us to choose $D$ before
taking any limit in $q$.

We first compare the probability cost of a point profile with the sum
of its local gains.  For sufficiently large $D$,
Lemma~\ref{lem:tail-residue} reduces the remaining cases to zero common
dimensions and pair dimensions $0$ or $2$.
Lemma~\ref{lem:tail-residual-counts} then gives stronger counts for these
cases, separating pair spaces that form a direct sum from those that do not.

\subsection{Exposing a profile}

Recall that $a_e$ is the fraction of the three bit positions at which the
two faces on $e$ differ.  Proposition~\ref{prop:complex} gives
$a_e\ge 1/3$ for every $e$, and $a_e>1/3$ for at least one $e$.

\begin{proposition}[Global profile cost]\label{prop:global-cost}
For every prescribed profile, its probability is $O_D(q^{-\mathcal C})$,
where
\begin{equation}\label{eq:global-cost}
 \mathcal C=\frac12\left(
 \sum_{j\in J}C_j-\sum_{e\in E}(1-a_e)b(h_e)\right).
\end{equation}
\end{proposition}

\begin{proof}
Fix one of the six classes obtained by prescribing one bit position and
its value.  It consists of eleven faces, is connected in the face
adjacency graph, and uses all thirteen points.  Choose a spanning tree
and a parent-first order.  The first face uses three points.  Every
subsequent face shares a pair with its parent and can add at most one
point.  Since all thirteen points occur, the ten subsequent faces must
each add exactly one fresh point.

Expose all the point Lagrangians occurring in the first face.  Its
prescribed profile has probability $O_D(q^{-C_j})$ by
Proposition~\ref{prop:triple-orbits}.  At a subsequent face $j$, let $e$
be its pair shared with its parent, and condition on the sigma-field
generated by \emph{all} previously exposed point Lagrangians.  Earlier
face-profile events are measurable in this sigma-field.  The fresh
point remains independent and uniform.  If the already exposed pair
has the prescribed intersection dimension $h_e$, pair transitivity and
Proposition~\ref{prop:triple-orbits} give the conditional probability
\[
 \frac{\P(\text{the prescribed profile on }j)}
      {\P(\dim(L_i\cap L_k)=h_e)}
 =O_D(q^{-C_j+b(h_e)}).
\]
Here pair transitivity makes the ratio the same for every \emph{specific}
pair of that dimension.  Thus the bound holds after conditioning on
the full earlier configuration, however its older points are related.
If the pair has the wrong dimension, the conditional probability is
zero.  Iterated conditional expectation proves a class probability
bound with cost $\sum_{j\text{ in class}}C_j$ minus $b(h_e)$ for its tree
pairs.  Subtracting $b(h_e)$ for every internal pair of the class only
weakens this bound.

Let $A$ be the complete profile event and let $A_1,\ldots,A_6$ be the six
class profile events.  Because $A\subseteq A_t$ for each $t$,
\[
 \P(A)^6\le\prod_{t=1}^6\P(A_t).
\]
This inequality does not require independence among the classes, all
of which use all thirteen points.  Each face appears in three class
bounds.  A pair $e$ is internal in exactly $3(1-a_e)$ classes, one for
each position where its two faces have the same bit.  Taking sixth
roots gives Equation~\eqref{eq:global-cost}.
\end{proof}

If the union in Proposition~\ref{prop:planted-gain} is empty for any
three-active face, its factor is zero and the entire profile contribution
vanishes; omit that profile. For each remaining profile and each face,
choose an exponent $g_j$ with $F_j=O_D(q^{g_j})$ as in
Proposition~\ref{prop:planted-gain} and Lemma~\ref{lem:two-active-gain}.
For three active incidences choose a maximizing parameter tuple in
Equation~\eqref{eq:planted-gain}. The choices range over finite sets
at fixed $D$. The pointwise gain bounds
and Proposition~\ref{prop:global-cost} imply
\begin{equation}\label{eq:tail-profile-bound}
 \E\left[\1_{\text{profile}}\prod_jF_j\right]
 =O_D\left(q^{\sum_jg_j-\mathcal C}\right).
\end{equation}
In particular, center weighting introduces no new conditional
independence assertion in the preceding exposure argument.

\subsection{Local slack and its quadratic part}

The profile bound gives decay when its probability cost exceeds the
sum of its local gains. We separate this difference into face
contributions and a nonnegative term from bit separation. A quadratic
bound will control large intersection dimensions; separate span
counts will handle the equality profiles.

Define the baseline and slack on a face by
\[
 \beta_j=\frac{C_j}{2}-\frac16\sum_{e\subset j}b(h_e),
 \qquad S_j=\beta_j-g_j,
\]
and define the nonnegative excess separation term
\begin{equation}\label{eq:tail-excess}
 \mathcal E=\frac12\sum_{e\in E}(a_e-1/3)b(h_e).
\end{equation}
Since each pair occurs twice,
\begin{equation}\label{eq:tail-total-slack}
\mathcal C-\sum_jg_j=\sum_j S_j+\mathcal E.
\end{equation}

Local slack can be negative. On a fully active face with
\(L_1=L_2=L_3\), all common and pair dimensions equal \(D\), and
\(F_j=\pi_{D,r}^{-2}\). Its gain exponent is \(2b(r)\), whereas
\(\beta_j=b(D)/2\), so \(S_j=-D/4\). The global sum and separation
below control this common-intersection mode.

\begin{lemma}[Faces with at most two active incidences]\label{lem:tail-lower-active}
Let $h=\max_{e\subset j}h_e$.  If $j$ has at most two active incidences,
then
\begin{equation}\label{eq:tail-lower-active}
 S_j\ge0,
 \qquad S_j\ge\frac{h^2}{24}-\frac h{12}.
\end{equation}
Equality $S_j=0$ requires $c_j=0$ and either all three pair dimensions
zero or exactly one pair dimension equal to $2$.  The latter case
requires exactly two active incidences, at that pair's endpoints.
\end{lemma}

\begin{proof}
Put $m=\max_\ell s_\ell$ and choose $\ell_0$ with $s_{\ell_0}=m$.
Using $h=c+m$ and $\sum s_\ell\le D-c$, expand
\begin{align*}
 3C-\sum_\ell b(c+s_\ell)-2b(c+m)
 &=3(D+1)c-5b(c)
   -c\left(\sum_\ell s_\ell+2m\right)
   +2\sum_{\ell\ne\ell_0}b(s_\ell)\\
 &\ge b(c)+2\sum_{\ell\ne\ell_0}b(s_\ell).
\end{align*}
The last step uses $\sum s_\ell+2m\le3(D-c)$; substituting this upper
bound leaves exactly $b(c)$ in the terms involving $c$.  Consequently
\begin{equation}\label{eq:tail-baseline-lower}
 \beta_j\ge\frac{b(h)}3+\frac{b(c)}6
             +\frac13\sum_{\ell\ne\ell_0}b(s_\ell).
\end{equation}
The gain is zero with zero or one active incidence.  With two, it is
$b(\lfloor h'/2\rfloor)$ for their pair dimension $h'\le h$.  For an
integer $h\ge0$,
\[
 b(\lfloor h/2\rfloor)\le\frac{b(h)}3,
 \qquad b(\lfloor h/2\rfloor)\le\frac{h(h+2)}8.
\]
The first follows separately for $h=2k$ and $h=2k+1$; its differences
are $k(k-1)/6$ and $(k+1)(k+2)/6$, respectively, so equality holds
exactly for $h=0,2$.  Combining these inequalities with
Equation~\eqref{eq:tail-baseline-lower} proves nonnegativity and
$S_j\ge b(h)/3-h(h+2)/8=h^2/24-h/12$.
Equality forces $c=0$, $s_\ell=0$ for $\ell\ne\ell_0$, and $h=0$ or
$2$.  When $h=2$, the gain must equal $1$, which requires the stated
two active endpoints.  With at most one active incidence the positive
baseline for $h>0$ excludes equality.
\end{proof}

For a face with three active incidences, retain the notation of
Section~\ref{sec:planted}:
\[
 \begin{gathered}
 u,\quad r'=r-u,\quad\delta=2u-c=d-2r',\quad
 s_\ell=p_\ell+t_\ell,\\
 P=\sum_\ell p_\ell,\quad T=\sum_\ell t_\ell,
 \quad w=\frac{d-3\delta}{2}-P+v.
 \end{gathered}
\]
The necessary constraints, which will suffice below, are
\begin{equation}\label{eq:tail-feasible}
 \begin{gathered}
 c\ge|\delta|,\qquad P+T\le d,\qquad
 p_\ell,t_\ell,v\ge0,\qquad
 t_\ell,v\le p_\ell+\delta,\qquad w\ge0.
 \end{gathered}
\end{equation}
In particular $T\le P+3\delta$, and hence
$d+3\delta\ge P+T+3\delta\ge2T\ge0$.
All the dimensions are integers, and $\delta=2u-c$ with $0\le u\le c$.

Substituting the gain in Equation~\eqref{eq:planted-gain} into
$S=\beta-g$ gives the exact identity
\begin{equation}\label{eq:slack}
 \begin{split}
 S={}&\frac{c(d+\delta-1)-\delta(d+2)+2\delta^2}{4}\\
 &+\sum_{\ell=1}^3\left[
 \frac{(2p_\ell-t_\ell)(2p_\ell-t_\ell-1)}6
 +\delta p_\ell-\frac c6(p_\ell+t_\ell)\right]
 +vw-b(v).
 \end{split}
\end{equation}
For clarity, the algebra separates into two identities.  The
common-space part uses $D=d+c$ and $u=(c+\delta)/2$:
\[
 \frac{(D+1)c}{2}-\frac{b(c)}2
 -\big((r+c+1)u-2u^2\big)
 =\frac{c(d+\delta-1)-\delta(d+2)+2\delta^2}{4}.
\]
For each reduced pair, use $b(c+s)=b(c)+b(s)+cs$ and
\[
 p^2+b(t)-\frac23b(p+t)
 =\frac{(2p-t)(2p-t-1)}6.
\]
These expansions yield every term of Equation~\eqref{eq:slack}.

Let \(Q\) be the homogeneous quadratic part of its right-hand side:
\begin{equation}\label{eq:tail-quadratic}
\begin{split}
 Q={}&\frac{c(d+\delta)-\delta d+2\delta^2}{4}\\
 &+\sum_{\ell=1}^3\left[
 \frac{(2p_\ell-t_\ell)^2}{6}
 +\delta p_\ell-\frac c6(p_\ell+t_\ell)\right]
 +v\left(\frac{d-3\delta}{2}-P\right)+\frac{v^2}{2}.
\end{split}
\end{equation}
The remaining terms are linear:
\begin{equation}\label{eq:tail-linear-loss}
 S=Q-\frac c4-\frac\delta2-\frac{2P-T}{6}-\frac v2.
\end{equation}

The next estimate will turn the total slack into bounds on
\(P,T,\delta\). When a pair dimension grows, these bounds force its
common part to spread through neighboring faces until the excess
separation is too large. Once all pair dimensions are bounded, we
return to the exact slack and use its coefficient of \(d\).

\begin{lemma}[Coercivity]\label{lem:coercivity}
Every tuple satisfying Equation~\eqref{eq:tail-feasible} satisfies
\begin{equation}\label{eq:coercivity}
 Q\ge\frac1{342}(P^2+T^2+\delta^2).
\end{equation}
\end{lemma}

\begin{proof}
Set
\[
 a=\left(P-\frac{d-3\delta}{2}\right)_+,
 \qquad v_0=\min\left(\frac P3+\delta,
                         \frac{P-T+3\delta}{2}\right).
\]
Both arguments of the minimum are nonnegative, by
Equation~\eqref{eq:tail-feasible}.  Completing the square gives
\[
 vw-\frac{v^2}{2}
 =\left(\frac{d-3\delta}{2}-P\right)v+\frac{v^2}{2}
 \ge-\frac{a^2}{2}.
\]
The condition $w\ge0$ gives $a\le v$; furthermore
$v\le\min_\ell(p_\ell+\delta)\le P/3+\delta$.  From $d\ge P+T$ we
also obtain $a\le(P-T+3\delta)/2$.  Therefore $a\le v_0$, and
\begin{equation}\label{eq:tail-scalar-quadratic}
 \begin{split}
 Q\ge{}&\frac{c(d+\delta)-\delta d+2\delta^2}{4}
 +\frac{(2P-T)^2}{18}+\delta P
 -\frac{c(P+T)}6-\frac{v_0^2}{2}.
 \end{split}
\end{equation}
Here we also used $\sum_\ell(2p_\ell-t_\ell)^2\ge(2P-T)^2/3$.
The coefficient of $c$ on the right is
\[
 \frac{d+\delta}{4}-\frac{P+T}{6}
 \ge\frac{d+3\delta}{12}\ge0.
\]
Because \(c\ge|\delta|\), replacing \(c\) by \(|\delta|\) can only
decrease the right side. We do so and treat the two signs of \(\delta\).

First suppose $\delta\ge0$.  The remainder in
Equation~\eqref{eq:tail-scalar-quadratic} becomes
\[
 F(T)=\frac{3\delta^2}{4}+\frac{(2P-T)^2}{18}
            +\frac{(5P-T)\delta}{6}-\frac12
            \min\left(\frac P3+\delta,
                        \frac{P-T+3\delta}{2}\right)^2,
 \qquad 0\le T\le P+3\delta.
\]
On $0\le T\le P/3+\delta$, the first argument attains the minimum,
and $F'(T)=(T-2P)/9-\delta/6\le0$; at the right endpoint this
derivative is $-5P/27-\delta/18$.  On the remaining interval the
coefficient of $T^2$ is $1/18-1/8=-5/72$, so $F$ is concave.  Its
minimum is therefore at one of the following two endpoints:
\begin{align*}
 F(P/3+\delta)&=\frac{32P^2+84P\delta+45\delta^2}{324},\\
 F(P+3\delta)&=\frac{2P^2+12P\delta+27\delta^2}{36}.
\end{align*}
Both are at least $(P^2+\delta^2)/18$.  Moreover
\[
 P^2+T^2+\delta^2
 \le P^2+(P+3\delta)^2+\delta^2
 \le3P^2+19\delta^2\le19(P^2+\delta^2).
\]
This proves the desired coefficient $1/(18\cdot19)$ in this case,
including $P=\delta=0$.

Next suppose $\delta=-y<0$.  Put $x=P-3y\ge0$; then $0\le T\le x$.
After setting $c=y$, the coefficient of $d$ is $y/2$.
Using $d\ge P+T=x+3y+T$ in
Equation~\eqref{eq:tail-scalar-quadratic} leaves
\[
 G(T)=\frac{y^2}{4}+\frac{(2x-T)y}{3}
 +\frac{(2x-T)^2}{18}
 -\frac12\min\left(\frac x3,\frac{x-T}{2}\right)^2.
\]
Its derivative up to $T=x/3$ is $(T-2x)/9-y/3\le0$; after that
breakpoint its quadratic coefficient is again $-5/72$.  Thus its
minimum is at least the smaller of its breakpoint and last-endpoint values,
\begin{align*}
 G(x/3)&=\frac{32x^2+180xy+81y^2}{324},\\
 G(x)&=\frac{2x^2+12xy+9y^2}{36}.
\end{align*}
Each is at least $(x^2+y^2)/18$, including the degenerate interval
$x=0$.  Finally,
\[
 P^2+T^2+y^2\le(x+3y)^2+x^2+y^2
                 \le3x^2+19y^2\le19(x^2+y^2).
\]
The same constant $1/342$ follows, completing the proof.
\end{proof}

\subsection{Uniform control and equality profiles}

Put $M=\max_{e\in E}h_e$.  On every face with three active incidences,
$c\le M$, $P\le3(M-c)$, $\delta\le c$, and
$v\le p_\ell+\delta\le s_\ell+c=h_\ell\le M$.
Equation~\eqref{eq:tail-linear-loss} consequently gives the uniform bound
\begin{equation}\label{eq:tail-uniform-linear}
 S\ge Q-\frac{3M}{2}.
\end{equation}
Indeed the subtracted expression is at most
$c/4+c/2+(M-c)+M/2=3M/2-c/4$.
Suppose that the cost in Equation~\eqref{eq:global-cost} does not yet
give a strict saving, so $\sum_jS_j+\mathcal E\le0$.
Let $J_3$ denote the faces with three active incidences and let
$h_j=\max_{e\subset j}h_e$.  Summing
Lemma~\ref{lem:coercivity}, Equation~\eqref{eq:tail-uniform-linear},
and Lemma~\ref{lem:tail-lower-active} gives
\begin{equation}\label{eq:tail-budget}
 \begin{split}
 &\sum_{j\in J_3}
 \frac{P_j^2+T_j^2+\delta_j^2}{342}
 +\sum_{j\notin J_3}\frac{h_j^2}{24}+\mathcal E
 \le33M.
 \end{split}
\end{equation}
To see the constant explicitly, each three-active face can lose at
most $3M/2$, and each lower-activity face loses at
most $h_j/12\le M/12\le3M/2$; there are twenty-two faces in total.
Every term on the left of
Equation~\eqref{eq:tail-budget} is nonnegative.  Thus it bounds each
term separately, despite the possible negative individual slacks.

\begin{lemma}[Profiles not separated by the first cost]\label{lem:tail-residue}
There is an even-dimensional threshold such that, for every even $D$
above it, every feasible profile and every activity set for which
$\sum_jg_j\ge\mathcal C$ has the following properties:
\begin{equation}\label{eq:tail-residual-profile}
 c_j=0\quad(j\in J),\qquad h_e\in\{0,2\}\quad(e\in E),
 \qquad\sum_jg_j\le\frac23\sum_{e\in E}b(h_e).
\end{equation}
\end{lemma}

\begin{proof}
Assume \(\sum_j S_j+\mathcal E\le0\). We first bound all pair
dimensions independently of \(D\), the activity set, and the choices
of maximizing tuples.

\paragraph{A uniform bound for pair dimensions.}
Equation~\eqref{eq:tail-budget} gives a constant \(K\ge1\), depending
only on the fixed numerical coefficients, such that
\[
 P_j+T_j\le K\sqrt M\quad(j\in J_3),\qquad
 h_j\le K\sqrt M\quad(j\notin J_3),\qquad
 \mathcal E\le33M.
\]
The case \(M=0\) already has the required bound. Suppose \(M>0\),
and choose a face \(j_0\) containing a pair of dimension \(M\).
For sufficiently large \(M\), this face cannot have lower activity,
because \(M>K\sqrt M\). On a three-active face every reduced pair
dimension satisfies \(s_\ell\le P+T\le K\sqrt M\). Thus a pair of
dimension \(h\) on that face forces its common dimension, and hence
every pair dimension on the face, to be at least \(h-K\sqrt M\).

Let \(\ell_*\) be the diameter of the fixed connected face-neighbor
graph. We claim by induction on \(t\) that every face at distance
\(t\) from \(j_0\) is three-active and every pair on it has dimension
at least \(M-(t+1)K\sqrt M\). The preceding paragraph gives the case
\(t=0\). When \(M\) is large enough that
\[
 M-(\ell_*+1)K\sqrt M\ge M/2>K\sqrt M,
\]
choose for a face at distance \(t+1\) its predecessor on a shortest
path from \(j_0\). Their shared pair has dimension greater than
\(K\sqrt M\), so the new face cannot have lower activity. Applying
the three-active bound proves the claimed lower bound at distance
\(t+1\).
Induction therefore gives \(h_e\ge M/2\) for every pair \(e\).

Fix a pair \(e_*\) with \(a_{e_*}>1/3\). Its contribution to
\(\mathcal E\) is at least
\[
 \frac12(a_{e_*}-1/3)b(h_{e_*})
 \ge \frac{a_{e_*}-1/3}{16}M^2.
\]
For large \(M\) this contradicts \(\mathcal E\le33M\). Hence there is
an integer \(M_*\), independent of \(D\), the activity set, and the
maximizing tuples, such that every profile under consideration
satisfies \(M\le M_*\).

\paragraph{Choose the dimension after this bound.}
On a three-active face with \(M\le M_*\), the integers
\(c,p_\ell,t_\ell,v\) lie between \(0\) and \(M_*\), and
\(|\delta|\le M_*\). Put \(n_\ell=2p_\ell-t_\ell\). Substituting the
definition of \(w\) into Equation~\eqref{eq:slack} gives
\begin{equation}\label{eq:tail-leading-coefficient}
 S=Ad+R,\qquad
 A=\frac{c-\delta+2v}{4}=\frac{c-u+v}{2},
\end{equation}
where
\[
\begin{split}
 R={}&\frac{c(\delta-1)-2\delta+2\delta^2}{4}
 +\frac16\sum_{\ell=1}^3
 \left[n_\ell(n_\ell-1)+6\delta p_\ell-c(p_\ell+t_\ell)\right]\\
 &-\left(\frac{3\delta}{2}+P\right)v+\frac{v(v-1)}2.
\end{split}
\]
Let \(C_*\) be the maximum of \(|R|\) over the finite integer box
\(0\le c,p_\ell,t_\ell,v\le M_*\), \(|\delta|\le M_*\), with
\(P=\sum_\ell p_\ell\). This bounds every actual remainder and is
independent of \(D\).

For an actual tuple, \(0\le u\le c\) and \(v\ge0\), so \(A\ge0\).
Moreover \(c,u,v\) are integers: either \(A=0\) or \(A\ge1/2\).
Choose an even \(D\) with
\[
 D>M_*+44C_*.
\]
If a three-active face had \(A>0\), it would have
\(S\ge(D-M_*)/2-C_*\). Every other three-active face has
\(S\ge-C_*\), every lower-activity face has nonnegative slack, and
\(\mathcal E\ge0\). Since there are \(22\) faces, the total slack
would be at least
\[
 \frac{D-M_*}{2}-22C_*>0,
\]
contrary to the hypothesis. Thus every three-active face has
\(u=c\), \(\delta=c\), and \(v=0\).

\paragraph{Integer equality profiles.}
For these parameters, Equation~\eqref{eq:slack} becomes
\begin{equation}\label{eq:tail-integer-slack}
 S=\frac{3c^2-3c}{4}
    +\frac16\sum_{\ell=1}^3
       \left[n_\ell^2+(c-1)n_\ell+3cp_\ell\right].
\end{equation}
If \(c\ge2\), completing the square in \(n_\ell\) gives
\[
 S\ge\frac{3c(c-1)}4-\frac{(c-1)^2}{8}
       =\frac{(c-1)(5c+1)}8>0.
\]
If \(c=1\), then \(S=\sum_\ell(n_\ell^2+3p_\ell)/6\ge0\),
and equality forces \(p_\ell=t_\ell=0\) for every \(\ell\).
If \(c=0\), the constraint \(t_\ell\le p_\ell\) gives the integer
\(n_\ell\ge p_\ell\ge0\). Each \(n_\ell(n_\ell-1)\) is nonnegative,
and its zeros give precisely
\[
 (p_\ell,t_\ell)=(0,0)\quad\text{or}\quad(1,1).
\]
Thus every slack is nonnegative at the chosen dimension. The assumed
inequality forces every slack, and \(\mathcal E\), to be zero.

A zero-slack face with \(c=1\) has all pair dimensions equal to \(1\).
Any face across one of these pairs cannot be a lower-activity equality
case from Lemma~\ref{lem:tail-lower-active}, and cannot be a
three-active equality case with \(c=0\). It must have \(c=1\) as well.
Connectivity would give \(h_e=1\) for every \(e\), contradicting
\(\mathcal E=0\) at a pair with \(a_e>1/3\). The remaining equality
cases have \(c_j=0\) and \(h_e\in\{0,2\}\). As \(S_j=0\), their gains
equal their baselines, and
\[
 \sum_jg_j=\sum_j\frac13\sum_{e\subset j}b(h_e)
          =\frac23\sum_e b(h_e).
\]
The choices of \(M_*\) and \(C_*\) were uniform in the activity set
and maximizing tuples, so the dimension threshold works for all of
them simultaneously.
\end{proof}

\subsection{Two counts for the remaining pair spaces}

It remains to improve the cost for the profiles in
Equation~\eqref{eq:tail-residual-profile}. Their nonzero pair
intersections are two-dimensional, and each pair space must lie in
the Lagrangians at both endpoints. A dependence among the pair
spaces forces at least one extra dimension in the sum of the local
endpoint spans, producing a factor \(q^{-D}\). For a direct family,
the independent orthogonality constraints instead give the required
cost after the endpoint containment factors cancel. The total
ambient span need not be isotropic.

\begin{lemma}[A strict span inequality]\label{lem:tail-span}
Let $H_e$ be subspaces assigned to distinct pairs $e$ of a simple
graph.  Set
\[
 k=\dim\sum_e H_e,\qquad k_i=\dim\sum_{e\ni i}H_e.
\]
Then $\sum_i k_i\ge2k$.  If the nonzero spaces do not form a direct
sum, then $\sum_i k_i\ge2k+1$.
\end{lemma}

\begin{proof}
For any ordering of the spaces, let $\Delta_e$ be the increment in
the dimension of the global span when $H_e$ is added.  At each
endpoint $i$ of $e$, let $\Delta_{e,i}$ be its increment in the local
span.  The earlier local span is contained in the earlier global
span, so $\Delta_{e,i}\ge\Delta_e$.  Summing gives
$\sum_i k_i\ge2\sum_e\Delta_e=2k$.

If the family is not direct, choose an inclusion-minimal subfamily
$\mathcal T$ which is not direct.  Order its members first and let
$e$ be its last member.  Put
\[
 \delta_*=\dim\left(H_e\cap
          \sum_{f\in\mathcal T\setminus\{e\}}H_f\right)>0.
\]
At least one endpoint $i$ of $e$ fails to belong to some other member
of $\mathcal T$: otherwise every other pair in $\mathcal T$ would
contain both endpoints of $e$ and would equal $e$.  The members of
$\mathcal T$ visible at this $i$, including $e$, are therefore a
proper subfamily and form a direct sum by minimality.  Consequently
\[
 \Delta_{e,i}=\dim H_e=\Delta_e+\delta_*.
\]
All other increment comparisons remain nonnegative, and hence
$\sum_i k_i\ge2k+\delta_*\ge2k+1$.  The argument uses the actual
dimension $\delta_*$ and in particular includes a partial intersection
of $H_e$ with the earlier span.
\end{proof}

\begin{lemma}[Costs for the remaining profiles]\label{lem:tail-residual-counts}
For a profile satisfying Equation~\eqref{eq:tail-residual-profile},
let $m$ be the number of pairs with $h_e=2$.  The contribution of
configurations whose pair spaces are not direct is
$O_D(q^{-D+1641})$.  The contribution of configurations whose pair
spaces are direct is
\[
 O_D\left(q^{-\frac13\sum_e b(h_e)}\right).
\]
Here ``contribution'' includes the product of the center factors.
The numerical remainder $1641$ is independent of $D$.
\end{lemma}

\begin{proof}
Write $H_e=L_i\cap L_k$ for $e=\{i,k\}$, and use $k,k_i$ as in
Lemma~\ref{lem:tail-span}.  We count possible $H_e$ and impose their
containment in the independent $L_i$.  Discarding exactness of the
intersections can only enlarge the count.  A nonisotropic local span
cannot be contained in a Lagrangian and contributes zero.  For every
isotropic local span of dimension $k_i$, its containment probability
is $O_D(q^{-Dk_i+\binom{k_i}{2}})$ by
Lemma~\ref{lem:lagrangian-counts}.

For the nondirect case, first choose an arbitrary $k$-space in the
$2D$-space, in $O_D(q^{k(2D-k)})$ ways, and then choose the $m$
two-spaces inside it, in $O_D(q^{2m(k-2)})$ ways.  Requiring their
span to be the chosen space would only reduce this upper bound.  On
the stratum with fixed $k$ and $k_i$, multiplying the containment
probabilities gives exponent
\[
 -D\left(\sum_i k_i-2k\right)
 -k^2+2mk-4m+\sum_i\binom{k_i}{2}.
\]
The point degrees in the pair graph, given by
Proposition~\ref{prop:complex}, are
$(4,6,4,6,5,4,5,5,5,6,5,6,5)$.  Therefore
\[
 \sum_i\binom{k_i}{2}
 \le\sum_i\binom{2\deg(i)}2=618.
\]
Since $m\le33$ and $-k^2+2mk\le m^2$, the part of the exponent
independent of $D$ is at most
\[
 m^2-4m+618\le1575.
\]
Lemma~\ref{lem:tail-span} makes the coefficient of $-D$ at least
one.  There are only boundedly many $k,k_i$, independently of $D$,
because $k\le2m\le66$ and $k_i\le2\deg(i)$.  Thus the probability
of this part of the profile is $O_D(q^{-D+1575})$.
Its total gain is at most $\tfrac23\sum_e b(h_e)=2m\le66$, proving
the first bound.  The ambient span in this argument was arbitrary;
only the local spans were required to be isotropic.

For the direct case, the following count applies to arbitrary fixed
dimensions $h_e$.  Order the pairs and choose an ordered basis of
each $H_e$.  Each new basis vector must be orthogonal to the previous
vectors of its own space and to all vectors of earlier spaces sharing
an endpoint.  In a direct family these previous vectors are linearly
independent.  Nondegeneracy of the symplectic form identifies ambient
vectors with linear functionals, so their orthogonality equations are
independent, even if the previous vectors are mutually orthogonal.
If there are $a$ such equations, there are exactly $q^{2D-a}$ vectors
satisfying them.  Excluding choices which would destroy directness
only decreases this count.

Let $e\sim f$ mean distinct pairs sharing an endpoint, and count each
unordered pair $\{e,f\}$ once.  Dividing the sequential basis count
by the number of ordered bases in each space gives at most
\[
 O_D\left(q^{\,2D\sum_eh_e-\sum_eh_e^2
 -\sum_e\binom{h_e}{2}-\sum_{e\sim f}h_eh_f}\right)
\]
possible direct families satisfying all local orthogonality conditions.
Here $k_i=\sum_{e\ni i}h_e$, so
\[
 \sum_i k_i=2\sum_eh_e,
 \qquad
 \sum_i\binom{k_i}{2}
 =2\sum_e\binom{h_e}{2}+\sum_{e\sim f}h_eh_f.
\]
Distinct pairs in a simple graph share at most one endpoint, which
explains the coefficient of the last sum.  Multiplying by the
containment probabilities cancels both the $D$ terms and the cross
terms.  The remaining exponent is exactly
\[
 -\sum_eh_e^2+\sum_e\binom{h_e}{2}=-\sum_e b(h_e).
\]
The gain bound in Equation~\eqref{eq:tail-residual-profile} now leaves
$-\tfrac13\sum_e b(h_e)$, as asserted.
\end{proof}

\subsection{Completion of the rank-layer estimate}

\begin{proposition}[Vanishing singular contribution]\label{prop:singular-tail}
For every sufficiently large even $D=2r$ and every activity set
$B\subseteq\mathscr C$,
\begin{equation}\label{eq:tail-vanishing}
 \E_{L,Y}\left[
 \1_{\{\exists e=\{i,k\}\in E:\,\dim(L_i\cap L_k)>0\}}
 \prod_{(i,j)\in B}J(L_i,Y_j)\right]=o(1)
\end{equation}
as $q\to\infty$ through odd primes, with $D$ fixed.  The dimension
threshold can be chosen for all activity sets simultaneously.
\end{proposition}

\begin{proof}
Choose an even $D$ above the threshold in
Lemma~\ref{lem:tail-residue} and also with $D>1641$.  These thresholds
refer only to dimensions and the fixed complex.  For any profile with
some $h_e>0$, either $\mathcal C>\sum_jg_j$, or it satisfies
Equation~\eqref{eq:tail-residual-profile}.  In the first case,
Equation~\eqref{eq:tail-profile-bound} gives a strictly negative
$q$-exponent.  In the second, the two bounds in
Lemma~\ref{lem:tail-residual-counts} are strictly negative because
$D>1641$ and $\sum_e b(h_e)>0$.

At this fixed \(D\), every selected gain is an integer and
\(6\mathcal C\) is an integer, because \(3a_e\) is an integer.
A strict first-cost saving is therefore at least \(1/6\). A singular
direct residual profile saves at least \(1\), and the nondirect case
also saves at least \(1\) because even \(D>1641\) implies \(D\ge1642\).
There are finitely many profiles and activity sets, and the gain
exponents were selected from the fixed finite unions in
Proposition~\ref{prop:planted-gain}. Summing their vanishing
contributions proves Equation~\eqref{eq:tail-vanishing}.
The local maximizing tuples need not be realizable simultaneously:
the residue argument uses only their necessary dimension inequalities
and the shared profile dimensions.
In particular, no simultaneous limit in $D$ and $q$ has been used.
The bound $D>1641$ is needed here only for the nondirect residual
case and is not asserted to be a threshold for the whole proof.
\end{proof}

\begin{proof}[Proof of Lemma~\ref{lem:rank-layer}]
Choose an even $D$ large enough for both
Proposition~\ref{prop:transverse-limit} and
Proposition~\ref{prop:singular-tail}.  For the matrix variables let
\[
 \mathcal T=\{\det(X_i-X_k)\ne0\text{ for every }\{i,k\}\in E\}.
\]
On $\mathcal T$, Proposition~\ref{prop:transverse-limit} gives the
right-hand side of Equation~\eqref{eq:rank-layer}, up to $o(1)$,
for every fixed activity set and every choice of incidence signs.
The integrand is nonnegative.  Lemma~\ref{lem:chart-comparison}
bounds its contribution on $\mathcal T^c$ by a constant depending
on $D$ times the unsigned Lagrangian contribution in
Equation~\eqref{eq:tail-vanishing}: graph intersections project isomorphically
onto the kernels of matrix differences, and dropping the prescribed signs
only enlarges the integrand after the comparable normalizations.
This contribution is $o(1)$ by
Proposition~\ref{prop:singular-tail}.  Adding the transverse and
singular parts proves Equation~\eqref{eq:rank-layer}.
There are only finitely many activity sets and sign assignments, so
the convergence is uniform over the choices required in
Section~\ref{sec:kernel}.  This proves Lemma~\ref{lem:rank-layer}.
The construction in Section~\ref{sec:kernel} therefore yields the strict
symmetric-kernel inequality, and Lemma~\ref{lem:finite-host} gives a finite
simple host $G$ with $t(H,G)<p(G)^{66}$, proving Theorem~\ref{thm:main}.
\end{proof}
